\section{Introduction}

Suppose Alice chooses a positive number $n$, Bob supplies a vector of length
$n$, and Carol selects a valid position in that vector. The interaction has
three owners and a genuine dependency: Alice's choice changes the types of
Bob's and Carol's choices. Each participant has a different interface to the
same event. Alice supplies $n$; Bob and Carol receive it. At the next event,
Bob becomes the supplier. No participant is permanently the verifier, the
environment, or the opponent.

A role-owned quantifier records this structure directly:
\[
 \Q_A(n:\Nat_{>0}).\;
 \Q_B(v:\mathsf{Vec}(\Bool,n)).\;
 \Q_C(j:\mathsf{Fin}(n)).\;G(n,v,j).
 \tag{1}
\]
Viewed locally, an owned choice has the shape of a dependent sum: supply a
witness and continue at its instance. A received choice has the shape of a
dependent product: accept the witness and continue at that instance. Globally,
there is one binder and one chosen value, shared by all the continuations.
This shared binding is the object we study.

Player-indexed choices have a well-developed semantic foundation.
Abramsky's multi-agent logic gives choices, quantifiers, role permutations,
and information constraints an intrinsic game interpretation
\cite{Abramsky2007}. Multirole logic supplies role-indexed proof rules and
multiparty cut \cite{XiWu2023}. Dependent session types make communicated
values available in continuation types \cite{ToninhoCairesPfenning2011}, and
Wu and Xi's multiparty system includes quantification interpreted separately
at each endpoint \cite{WuXi2018}. We use these constructions as the basis for
one explicit account of binding, dependent substitution, and disclosure.

The computational question is immediate: what happens when the endpoints
meet? In (1), Alice's term provides a number $a$. Every endpoint must then
advance to the continuation for that same $a$. The substitution changes
continuation types as well as programs. It is therefore insufficient to say
only that all participants follow the same sequence of messages. Their
dependent interfaces must agree on the values already selected.

There is a second question when information is private. In a sealed-bid
interaction, Alice may fix a value before Bob acts and disclose it afterward.
The two moments have different meanings. Fixing the value determines which
choice was made; publication determines who may use it. If Alice could choose
the value at the later moment, she could make it depend on Bob's intervening
action. A calculus that expresses commitment must preserve the identity of the
earlier choice across the disclosure boundary.

These questions lead to three design decisions.
First, ownership is part of a global binder rather than a permanent polarity
assigned to a participant. Second, values used in dependent types are stable
data, and communication acts by ordinary dependent substitution. Third,
availability is recorded separately from ownership. A participant may supply a
value that only some participants can presently observe.

\paragraph{Interaction and objectives.}
A script describes how participants act together. Its objectives describe how
the resulting execution is evaluated. These are separate pieces of
information. The same script may carry incompatible goals, a common goal,
numeric utilities, or a family of proof obligations. Its formation asserts
none of those goals. A typed endpoint must provide the results required by its
interface; a payoff attached to an execution remains an observation of the
interaction. This distinction lets the calculus retain the whole multi-role
object, including behaviors outside any selected solution concept.

\paragraph{The core developed here.}
The public calculus consists of finite dependent scripts, endpoint terms, and
their complete connections. Its data language is a fixed dependent type theory.
We give a sum/product interpretation and prove the substitution, interaction,
and role-renaming laws. Complete executions can be passed through typed public
return boundaries, giving an associative sequencing operation. We then give an
observation-based account of private binding and publication. The disclosure
results concern its specified transition semantics; a dependent typing system
for all information-changing networks is a further extension.

\paragraph{From a core to a larger calculus.}
The next composition problem is dependent multiparty cut across open
interfaces. A component may consume some endpoints and provide others; its
substitution must preserve both the dependent data and the permitted
observations. The final sections specify that problem and its proof
obligations. The distinction between complete scripted connections and general
open networks gives the development a concrete mathematical boundary.
